% GENERATED FILE. Edit canonical lessons, then run npm run generate:books.
% canonical-source-hash: fnv1a64:76ef88dfcfa3609e
% canonical-lessons: TE-C210-anumanam, TE-C210-trpti, TE-C210-nirasa, TE-C210-visugu, TE-C210-benga

\chapter{Doubt, Contentment, Disappointment, Boredom, Pining}
\label{ch:te-doubt-contentment-disappointment-boredom-pining}
\hlchaptermodality{210}

\begin{chapteropening}
\textbf{By the end of this chapter:} \emph{I can say doubt, contentment, disappointment, boredom and pining.}

Complete the last lesson of chapter 210: I can say doubt, contentment, disappointment, boredom and pining.
\end{chapteropening}

\section[\texorpdfstring{\te{అనుమానం} (anumānaṁ)}{అనుమానం (anumānaṁ)}]{\te{అనుమానం} (anumānaṁ) — doubt, suspicion}

\label{lesson:TE-C210-anumanam}



\begin{quote}
Before the new one: say the Telugu for medical treatment, then the Telugu for the forehead.
\end{quote}

\subsection*{You'll want to know: \te{అనుమానం}}
\textbf{\te{అనుమానం}} — \emph{anumānaṁ} — \textquotedblleft{}doubt, suspicion\textquotedblright{}.

\begin{grammarlens}[title={What you've built}]
One more word for this chapter.
\end{grammarlens}

\subsection*{Your turn}
\begin{itemize}
\raggedright
  \item \emph{Say it:} \emph{anumānaṁ}
  \item \emph{Say it:} \emph{anumānaṁ}, once more
  \item \emph{Say it:} say \emph{anumānaṁ}
  \item \emph{Recall:} say the Telugu for medical treatment, then the Telugu for the forehead, then say \emph{anumānaṁ} again
\end{itemize}

\begin{tcolorbox}[breakable,colback=teal!4,colframe=teal!35!black,title={Before you move on}]
What does \te{అనుమానం} mean? (\textquotedblleft{}Doubt, suspicion\textquotedblright{}.) Say it once more.
\end{tcolorbox}
\section[\texorpdfstring{\te{తృప్తి} (tṛpti)}{తృప్తి (tṛpti)}]{\te{తృప్తి} (tṛpti) — contentment, satisfaction}

\label{lesson:TE-C210-trpti}



\begin{quote}
Before the new one: say the Telugu for the forehead, then the Telugu for doubt.
\end{quote}

\subsection*{You'll want to know: \te{తృప్తి}}
\textbf{\te{తృప్తి}} — \emph{tṛpti} — \textquotedblleft{}contentment, satisfaction\textquotedblright{}.

\begin{grammarlens}[title={What you've built}]
One more word for this chapter.
\end{grammarlens}

\subsection*{Your turn}
\begin{itemize}
\raggedright
  \item \emph{Say it:} \emph{tṛpti}
  \item \emph{Say it:} \emph{tṛpti}, once more
  \item \emph{Say it:} say \emph{tṛpti}
  \item \emph{Recall:} say the Telugu for the forehead, then the Telugu for doubt, then say \emph{tṛpti} again
\end{itemize}

\begin{tcolorbox}[breakable,colback=teal!4,colframe=teal!35!black,title={Before you move on}]
What does \te{తృప్తి} mean? (\textquotedblleft{}Contentment, satisfaction\textquotedblright{}.) Say it once more.
\end{tcolorbox}
\section[\texorpdfstring{\te{నిరాశ} (nirāśa)}{నిరాశ (nirāśa)}]{\te{నిరాశ} (nirāśa) — disappointment}

\label{lesson:TE-C210-nirasa}



\begin{quote}
Before the new one: say the Telugu for doubt, then the Telugu for contentment.
\end{quote}

\subsection*{You'll want to know: \te{నిరాశ}}
\textbf{\te{నిరాశ}} — \emph{nirāśa} — \textquotedblleft{}disappointment\textquotedblright{}.

\begin{grammarlens}[title={What you've built}]
One more word for this chapter.
\end{grammarlens}

\subsection*{Your turn}
\begin{itemize}
\raggedright
  \item \emph{Say it:} \emph{nirāśa}
  \item \emph{Say it:} \emph{nirāśa}, once more
  \item \emph{Say it:} say \emph{nirāśa}
  \item \emph{Recall:} say the Telugu for doubt, then the Telugu for contentment, then say \emph{nirāśa} again
\end{itemize}

\begin{tcolorbox}[breakable,colback=teal!4,colframe=teal!35!black,title={Before you move on}]
What does \te{నిరాశ} mean? (\textquotedblleft{}Disappointment\textquotedblright{}.) Say it once more.
\end{tcolorbox}
\section[\texorpdfstring{\te{విసుగు} (visugu)}{విసుగు (visugu)}]{\te{విసుగు} (visugu) — boredom, weariness of something}

\label{lesson:TE-C210-visugu}



\begin{quote}
Before the new one: say the Telugu for contentment, then the Telugu for disappointment.
\end{quote}

\subsection*{You'll want to know: \te{విసుగు}}
\textbf{\te{విసుగు}} — \emph{visugu} — \textquotedblleft{}boredom, weariness of something\textquotedblright{}.

\begin{grammarlens}[title={What you've built}]
One more word for this chapter.
\end{grammarlens}

\subsection*{Your turn}
\begin{itemize}
\raggedright
  \item \emph{Say it:} \emph{visugu}
  \item \emph{Say it:} \emph{visugu}, once more
  \item \emph{Say it:} say \emph{visugu}
  \item \emph{Recall:} say the Telugu for contentment, then the Telugu for disappointment, then say \emph{visugu} again
\end{itemize}

\begin{tcolorbox}[breakable,colback=teal!4,colframe=teal!35!black,title={Before you move on}]
What does \te{విసుగు} mean? (\textquotedblleft{}Boredom, weariness of something\textquotedblright{}.) Say it once more.
\end{tcolorbox}
\section[\texorpdfstring{\te{బెంగ} (beṅga)}{బెంగ (beṅga)}]{\te{బెంగ} (beṅga) — pining, homesickness}

\label{lesson:TE-C210-benga}



\begin{quote}
Before the new one: say the Telugu for disappointment, then the Telugu for boredom.
\end{quote}

\subsection*{You'll want to know: \te{బెంగ}}
\textbf{\te{బెంగ}} — \emph{beṅga} — \textquotedblleft{}pining, homesickness\textquotedblright{}.

\begin{grammarlens}[title={What you've built}]
One more word for this chapter.
\end{grammarlens}

\subsection*{Your turn}
\begin{itemize}
\raggedright
  \item \emph{Say it:} \emph{beṅga}
  \item \emph{Say it:} \emph{beṅga}, once more
  \item \emph{Say it:} say \emph{beṅga}
  \item \emph{Recall:} say the Telugu for disappointment, then the Telugu for boredom, then say \emph{beṅga} again
\end{itemize}

\begin{tcolorbox}[breakable,colback=teal!4,colframe=teal!35!black,title={Before you move on}]
What does \te{బెంగ} mean? (\textquotedblleft{}Pining, homesickness\textquotedblright{}.) Say it once more.
\end{tcolorbox}
